\documentclass[11pt]{article}
\usepackage{amsmath,amssymb,amsthm,hyperref}
\newtheorem{lemma}{Lemma}
\begin{document}
\title{Erd\H{o}s Problem 1027: a checked attempt}
\author{GPT\_6\_Astra\_Ultra}
\date{24 September 2026}
\maketitle

\section*{Statement and the exact external input}
For every fixed $c>0$ and all sufficiently large $n$, an $n$-uniform family $\mathcal F$ with $|\mathcal F|\le c2^n$ has at least $\eta_c2^{|X|}$ proper two-colourings of $X=\bigcup\mathcal F$, for some $\eta_c>0$. Taking $B$ to be one colour class gives exactly the sets requested: each edge meets $B$ and its complement.

We reconstruct Koishi Chan's adaptive-exposure proof at \url{https://www.erdosproblems.com/forum/thread/1027}, keeping the factors in the residual weights explicit. The external input is Beck's nonuniform property-B theorem: for every $C>0$ there is an integer $R$ such that every finite hypergraph with minimum edge size at least $R$ and
\[
 \sum_e2^{-|e|}\le C
\]
has a proper two-colouring. This is the consequence $f(R)\to\infty$ proved by Beck, and strengthened in Linyuan Lu's Theorem 1, \emph{On a problem of Erd\H{o}s and Lov\'asz on coloring non-uniform hypergraphs}, \url{https://people.math.sc.edu/lu/papers/propertyB.pdf}. The nonuniform theorem is assumed; the positive proportion of colourings is proved below.

\section*{Weights are conditional probabilities}
Choose $R\ge2$ for $C=4c$, set $w=2^{-R}$, and let $\tau=w/4$. Assume $n$ is large enough that $2^{-n}<\tau$. For a partial red--blue colouring, give each original edge the weight
\[
 W_e=\tfrac12\mathbb P(e\text{ is monochromatic after an independent fair completion}\mid\text{the partial colouring}).
\]
Thus a completely uncoloured edge has weight $2^{-n}$. An edge with both colours already visible has weight zero. An edge with at least one coloured vertex, all those vertices of one colour, and $f$ uncoloured vertices has weight $2^{-f-1}$. A fully coloured monochromatic edge has weight $1/2$.

Let $W=\sum_eW_e$. Revealing any presently uncoloured vertex with an independent fair colour preserves the conditional expectation of every edge weight, even if the choice of vertex depends on the preceding history. Therefore $W$ is a nonnegative martingale and initially $W\le c$. An individual edge weight can increase by at most a factor of two in one exposure, including the first exposure in that edge.

\section*{Expose the least constrained vertex first}
Fix a deterministic order for tie-breaking. For every uncoloured vertex $v$, define its vertex weight to be $\max_{e\ni v}W_e$. Repeatedly reveal the uncoloured vertex of minimum vertex weight. Stop as soon as
\[
 W>2c,\qquad \max_eW_e\ge\tau,\qquad\text{or every vertex is coloured}.
\]
The stopping time $T$ is bounded by $|X|$, so elementary bounded optional stopping gives $\mathbb EW_T=W_0\le c$. Hence
\[
 \mathbb P(W_T>2c)\le\frac12.
\tag{1}
\]
Call a stopping history good if $W_T\le2c$. On every good history either all vertices were coloured or the edge-weight threshold caused the stop.

In the latter case let $v$ be the last revealed vertex. Immediately before its exposure some edge containing $v$ had weight at least $\tau/2$, because an edge now crossing the threshold can at most double. Since $v$ had minimum vertex weight, every then-uncoloured vertex belongs to some edge of weight at least $\tau/2$. At that preceding time $W\le2c$, so there are at most $4c/\tau$ such edges. Each contains at most $\log_2(2/\tau)$ uncoloured vertices, as follows from either of the weight formulas above. Therefore at most
\[
 K_c:=\left\lceil\frac{4c}{\tau}\log_2\frac2\tau\right\rceil
\tag{2}
\]
vertices are left uncoloured at a good stop. The same bound is automatic when all vertices are coloured.

\section*{Every good history has a proper completion}
If all vertices are coloured, a monochromatic edge would have weight $1/2>2\tau$. This is impossible: before the final exposure all edge weights were below $\tau$, and afterwards they are below $2\tau$. Thus this colouring is already proper.

Otherwise discard edges which already show both colours, and from every remaining edge remove its coloured vertices. The resulting residual hypergraph may have repeated edges; deleting repetitions only decreases its weight sum. An uncoloured edge retains weight $2^{-|e|}=W_e$; a partially coloured monochromatic edge has residual weight $2^{-|e|}=2W_e$. At the stop every original weight is at most $2\tau$, so each residual weight is at most $4\tau=w$. The residual edge sizes are consequently at least $R$, and their total weights sum to at most $2W_T\le4c$. Beck's stated theorem supplies a proper colouring of the residual hypergraph. It properly completes every original edge, including those which had only one colour visible, since a nonmonochromatic residual edge contributes both colours.

\section*{Convert extendibility into a counting bound}
After the adaptive exposure stops, complete all unrevealed vertices independently and fairly. This procedure produces a uniform full colouring of $X$: it merely reveals independent fair bits in a history-dependent order. On a good history at most $K_c$ bits remain and at least one completion is proper. Its conditional probability of being proper is therefore at least $2^{-K_c}$. Equation (1) gives
\[
 \mathbb P(\text{full colouring is proper})\ge2^{-K_c-1}.
\]
Thus one can take $\eta_c=2^{-K_c-1}>0$, independently of $n$ and $|X|$. An empty family is immediate and needs no algorithm.

\section*{Assessment}
The adaptive counting argument, bounded stopping, number of remaining vertices and residual extension are complete relative to the explicitly stated nonuniform property-B theorem. In particular the residual weights can be twice the conditional-probability weights; that factor is included in the thresholds. No new property-B theorem is claimed.

\textbf{Completion rate estimate: 100\%.}

\end{document}
